%%%%%%%%%%%%%%%%%%%%%%%%%%%%%%%%%%%%%%%%%%%%%%%%%%%%%%%%%%%%%%%%%%%%%%%%%%%%%%%
%
% Author: Marco Antonio Cardoso Moreno
% email: marcoacardosom@gmail.com
%
% Template para las tareas de la asignatura Diseño y Análisis de Algoritmos
%
%%%%%%%%%%%%%%%%%%%%%%%%%%%%%%%%%%%%%%%%%%%%%%%%%%%%%%%%%%%%%%%%%%%%%%%%%%%%%%%

% \documentclass[letterpaper, openright, 12pt]{book}
\documentclass[letterpaper, 12pt]{article}

\usepackage[spanish, es-tabla, es-nodecimaldot]{babel} % spanish
\usepackage[utf8]{inputenc} % for any non ASCII characters

\usepackage{graphicx}
\usepackage{caption}
\usepackage{subcaption}

\usepackage{physics} % simbolos matemáticos
\usepackage{amsmath} % mejora en fórmulas  matemáticas(aligns)
\usepackage{gensymb} % typesetting units (degrees, etc)
\usepackage{siunitx} % scientific notation and other stuff


\usepackage{listings} % for code blocks
\usepackage[left=2.5cm, top=2.5cm, right=2.5cm, bottom=2.5cm]{geometry}

\usepackage{pdfpages} % insert pdf as images
% \usepackage{cite} % referencias y bibliografía
% \usepackage{babelbib}

\usepackage{enumerate} %listas

% Cambiar el titulo por default de los indices
\addto{\captionsspanish}{\renewcommand*{\listfigurename}{Índice de Figuras}}
\addto{\captionsspanish}{\renewcommand*{\contentsname}{Índice}}
\addto{\captionsspanish}{\renewcommand*{\listtablename}{Índice de Tablas}}

% cambia nombre de Bibliografía a Referencias
% \renewcommand\bibname{Referencias}

% define estilo de cabeceras, número de páginas, etc
\pagestyle{plain}

% water mark images
% \usepackage[firstpageonly]{draftwatermark}
\usepackage{eso-pic}

% for phantom section* across the whole document. Better hyperrefs
\usepackage[hidelinks]{hyperref}
\providecommand\phantomsection*{}

% for references to equations, fiugres, etc
\usepackage{cleveref}

% new paragraph
\newcommand{\parg}{\par \hfill \break}
%%%%%%%%%%%%%%%%%%%%%%%%%%%%%%%%%%%%%%%%%%%%%%%%%%%%%%%%%%%%%%%%%%%%%%%%%%%%%%%
%
% Document
%
%%%%%%%%%%%%%%%%%%%%%%%%%%%%%%%%%%%%%%%%%%%%%%%%%%%%%%%%%%%%%%%%%%%%%%%%%%%%%%%

\begin{document}

% carátula
\begin{titlepage}
\begin{center}
  \AddToShipoutPictureBG*{
    \includegraphics[width=\paperwidth,height=\paperheight]
        {../../portada/portada.pdf}}
        \vspace*{2mm}
        \rule{\textwidth}{1.5pt}\\
        \vspace*{5mm}
        \textbf{\Large INSTITUTO POLITÉCNICO NACIONAL}\\
        \vspace*{8mm}
        {\Large Centro de Investigación en Computación}\\
        \vspace{8mm}
        {\Large Maestría en Ciencias en Ingeniería de Cómputo}\\
        \vspace{19mm}
        {\Large Diseño y Análisis de Algoritmos}\\
        \vspace{49mm}




        \textbf{\large Tarea 1 - Demostraciones de las propiedades de los órdenes asintóticos}\\





        \vspace*{48mm}
        {\large Docente: Dr.\ Rolando Quintero Téllez}\\
        \vspace*{20mm}
        {\large Alumno: Cardoso Moreno Marco Antonio}\\
        \vspace*{10mm}
        \rule{\textwidth}{1.5pt}\\
        {\large Semestre A21} \hfill {\large Fecha: 17-marzo-2021}
    \end{center}
\end{titlepage}

\newpage
\thispagestyle{empty}

% \newpage
% \pagenumbering{Roman}


% \clearpage
% \tableofcontents
% % \cleardoublepage
%
% \clearpage
% \addcontentsline{toc}{section*}{Índice de Figuras}
% \listoffigures
%
% \clearpage
% \addcontentsline{toc}{section*}{Índice de Tablas}
% \listoftables

\clearpage
\pagenumbering{arabic}

\section*{Demostrar:}
\label{sec:orgbe567fb}
\section{\(f(n)\) es \(\Theta(g(n))\) si y sólo si es \(O(g(n))\) y es \(\Omega(g(n))\)}
\label{sec:org5cd6a49}
\begin{itemize}
\item Si \(f(n)\) es \(\Theta(g(n))\), existen las constantes \(c_f', c_f\) y \(n_f\) tal
que
\begin{itemize}
\item \(c_f' g(n) \leq f(n) \leq c_f g(n)\) para toda \(n \geq n_f\).
\end{itemize}
\end{itemize}

\noindent
Si esta desigualdad la dividimos en dos, las nuevas desigualdades quedan:
$$
c_f' g(n) \leq f(n) \quad \mbox{y} \quad f(n) \leq c_f g(n)
    \quad \mbox{para toda} \quad n \geq n_f
$$

\noindent
Se observa que estas desigualdades son las definiciones tanto de
\(f(n) = O(g(n))\) (derecha), como de \(f(n) = \Omega(g(n))\) (izquierda), con lo
que se comprueba el enunciado.
\section{Transitividad}
\label{sec:org3ddb54e}
\subsection{Si \(f(n)\) es \(O(g(n))\) y \(g(n\)) es \(O(h(n))\) entonces \(f(n)\) es \(O(h(n))\)}
\label{sec:org350b076}
Como se comentó en los videos, debemos partir de que el antecedente es cierto,
por lo tanto, sabemos que:
\begin{itemize}
\item Si \(f(n)\) es \(O(g(n))\), existen dos constantes \(n_f\) y \(c_f\) tal que
\begin{itemize}
\item para toda \(n \geq n_f\), \(c_f g(n) \geq f(n)\)
\end{itemize}
\item Si \(g(n)\) es \(O(h(n))\), existen dos constantes \(n_g\) y \(c_g\) tal que
\begin{itemize}
\item para toda \(n \geq n_g\), \(c_g h(n) \geq g(n)\)
\end{itemize}
\end{itemize}

\noindent
Entonces:
$$
c_f c_g h(n) \geq f(n) \quad \mbox{para toda} \quad n \geq \max\{n_f, n_g\}
$$
\subsection{Si \(f(n)\) es \(\Theta(g(n))\) y \(g(n)\) es \(\Theta(h(n))\) entonces \(f(n)\) es \(\Theta(h(n))\)}
\label{sec:org588fd80}
Esta comprobación es una combinación de las pruebas de transitividad para \(O(n)\)
y \(\Omega(n)\). Asumiendo que el antecedente es cierto:
\begin{itemize}
\item Si \(f(n)\) es \(\Theta(g(n))\), existen las constantes \(c_f, c_f'\) y \(n_f\) tal que
\begin{itemize}
\item \(c_f' g(n) \leq f(n) \leq c_f g(n)\) para toda \(n \geq n_f\)
\end{itemize}
\item Si \(g(n)\) es \(\Theta(h(n))\), existen las constantes \(c_g, c_g'\) y \(n_g\) tal que
\begin{itemize}
\item \(c_g' h(n) \leq g(n) \leq c_g h(n)\) para toda \(n \geq n_g\)
\end{itemize}
\end{itemize}
\parg
Tenemos entonces:
$$
c_f' g(n) \leq f(n) \quad \mbox{y} \quad c_g' h(n) \leq g(n)
$$
si multiplicamos la ecuación de la derecha por la \textbf{constante positiva} \(c_f'\),
obtenemos:
$$
c_f' c_g' h(n) \leq c_f' g(n) \leq f(n) \quad \mbox{y por lo tanto} \quad
c_f' c_g' h(n) \leq f(n)
$$
\parg
\noindent
Por otro lado, tenemos:
$$
f(n) \leq  c_f g(n) \quad \mbox{y} \quad g(n) \leq c_g h(n)
$$
si multiplicamos la ecuación de la derecha por la \textbf{constante positiva} \(c_f\),
obtenemos:
$$
f(n) \leq c_f g(n) \leq c_f c_g h(n) \quad \mbox{y por lo tanto} \quad
f(n) \leq c_f c_g h(n)
$$
\parg

\noindent
Si unimos ambas partes, queda:
$$
c_f' c_g' h(n) \leq f(n) \leq c_f c_g h(n) \quad \mbox{para toda} \quad
n \geq \max\{n_f, n_g\}
$$
\subsection{Si \(f(n)\) es \(\Omega(g(n))\) y \(g(n)\) es \(\Omega(h(n))\) entonces \(f(n)\) es \(\Omega(h(n))\)}
\label{sec:org169c466}
Debemos asumir que el antecedente es cierto, entonces:
\begin{itemize}
\item Si \(f(n)\) es \(\Omega(g(n))\), existen dos constantes \(n_f'\) y \(c_f'\) tal que
\begin{itemize}
\item para toda \(n \geq n_f'\), \(c_f' g(n) \leq f(n)\)
\end{itemize}
\item Si \(g(n)\) es \(\Omega(h(n))\), existen dos constantes \(n_g'\) y \(c_g'\) tal que
\begin{itemize}
\item para toda \(n \geq n_g'\), \(c_g' h(n) \leq g(n)\)
\end{itemize}
\end{itemize}

Tenemos:
$$
c_f' g(n) \leq f(n)
$$
y
$$
c_g' h(n) \leq g(n)
$$
que al multiplicarse por la constante \textbf{positiva} \(c_f'\) queda
$$
c_f' c_g' h(n) \leq c_f' g(n)
$$
por lo que:
$$
c_f' c_g' h(n) \leq c_f' g(n) \leq f(n)
$$
entonces:

$$
c_f' f_g' h(n) \leq f(n) \quad \mbox{ para toda } \quad n \geq \max\{n_f', n_g'\}
$$
\section{Reflexividad}
\label{sec:orgcf4c494}
\subsection{Si \(f(n)\) es \(O(f(n))\) entonces \(f(n)\) es \(O(f(n))\)}
\label{sec:org57fdc4b}
\begin{itemize}
\item Si \(f(n)\) es \(O(f(n))\), existen dos constantes \(c_f\) y \(n_f\) tal que
\begin{itemize}
\item \(f(n) \leq c_f f(n)\) para toda \(n \geq n_f\)
\end{itemize}
\end{itemize}

\noindent
Si hacemos \(c_f = 1\), nos queda:
$$
f(n) \leq 1 \cdot f(n) \quad \mbox{para toda} \quad n \geq n_f
$$
o
$$
f(n) \leq f(n) \quad \mbox{para toda} \quad n \geq n_f
$$


\noindent
Con lo que se cumple la condición requerida, ya que \(f(n) = f(n)\).


\noindent
Si, por otro lado, despejamos \(c_f\) de \(f(n) \leq c_f f(n)\), tenemos:
\begin{equation*}
\frac{f(n)}{f(n)} \leq c_f \quad \mbox{por lo tanto} \quad 1 \leq c_f
\end{equation*}

de donde se observa que la condición se cumple para toda \(c_f \geq 1\).
\subsection{Si \(f(n)\) es \(\Omega(f(n))\) entonces \(f(n)\) es \(\Omega(f(n))\)}
\label{sec:org9cb5b2e}
\begin{itemize}
\item Si \(f(n)\) es \(\Omega(f(n))\), existen dos constantes \(c_f'\) y \(n_f'\) tal que
\begin{itemize}
\item \(c_f' f(n) \leq f(n)\) para toda \(n \geq n_f'\)
\end{itemize}
\end{itemize}

\noindent
Si de igual manera, hacemos \(c_f' = 1\) nos queda:
$$
1 \cdot f(n) \leq f(n) \quad \mbox{para toda} \quad n
$$
o
$$
f(n) \leq f(n) \quad \mbox{para toda} \quad n
$$

\noindent
Con lo que se cumple la condición requerida, ya que \(f(n) = f(n)\).

\noindent
Si, por otro lado, despejamos \(c_f'\) de \(c_f' f(n) \leq f(n)\), tenemos:
$$
c_f' \leq \frac{f(n)}{f(n)} \quad \mbox{por lo tanto} \quad c_f' \leq 1
$$
de donde se observa que la condición se cumple para toda \(c_f' \leq 1\). Es
importante recordar que en el análisis de los órdenes asintóticos las constantes
se asumen positivas, por lo que en este caso \(0 < c_f' \leq 1\).
\subsection{Si \(f(n)\) es \(\Theta(f(n))\) entonces \(f(n)\) es \(\Theta(f(n))\)}
\label{sec:orgdb0c2ff}
\begin{itemize}
\item Si \(f(n)\) es \(\Theta(f(n))\), existen las constantes \(c_f, c_f'\) y \(n_f\) tal
que \(c_f' f(n) \leq f(n) \leq c_f f(n)\)
\end{itemize}

\noindent
Si \(c_f = c_f' = 1\), tenemos:
$$
1 \cdot f(n) \leq f(n) \leq 1 \cdot f(n) \quad \mbox{para toda} \quad n
$$
o
$$
f(n) \leq f(n) \leq f(n) \quad \mbox{para toda} \quad n
$$
\section{Simetría}
\label{sec:org389def5}
\subsection{\(f(n)\) es \(\Theta(g(n))\) si y sólo si \(g(n)\) es \(\Theta(f(n))\)}
\label{sec:org792f820}
Tomando como cierto el antecedente, tenemos:
\begin{itemize}
\item Si \(f(n)\) es \(\Theta(g(n))\), existen las constantes \(c_f', c_f\) y \(n_f\) tal
que \(c_f' g(n) \leq f(n) \leq c_f g(n)\) para toda \(n \geq n_f\)
\end{itemize}

\noindent
Si a esta desigualdad, la separamos en dos
$$
c_f' g(n) \leq f(n) \quad \mbox{y} \quad f(n) \leq c_f g(n)
$$
si a su vez, a estas desigualdades las manipulamos algebráicamente de modo que
despejemos a \(g(n)\), nos quedan:
$$
g(n) \leq \frac{1}{c_f'} f(n) \quad \mbox{y} \quad \frac{1}{c_f} f(n) \leq g(n)
$$
de donde se infiere que \(\frac{1}{c_f'}\) y \(\frac{1}{c_f}\) son constantes
positivas, con lo que obtenemos
$$
\frac{1}{c_f} f(n) \leq g(n) \leq \frac{1}{c_f'} f(n)
$$

\noindent
Esto cumple los requerimientos para que \(g(n)\) sea \(O(f(n))\).
\section{Simetría Transpuesta}
\label{sec:orgc1f3492}
\subsection{\(f(n)\) es \(O(g(n))\) si y sólo si \(g(n)\) es \(\Omega(f(n))\)}
\label{sec:org4b4a0c3}
\begin{itemize}
\item Si \(f(n)\) es \(O(g(n))\), existen las constantes \(c_f\) y \(n_f\) tal que
\begin{itemize}
\item \(f(n) \leq c_f g(n)\) para toda \(n \geq n_f\)
\end{itemize}
\end{itemize}

\noindent
Si en dicha desigualdad despejamos \(g(n)\), tenemos:
$$
\frac{1}{c_f} f(n) \leq g(n) \quad \mbox{para toda} \quad n \geq n_f
$$
se observa entonces, por definición, que \(g(n)\) es \(\Omega(g(n))\).
\section{Aditividad}
\label{sec:orgc3c608e}
\subsection{Si \(f(n)\) es \(O(h(n))\) y \(g(n)\) es \(O(h(n))\), entonces \(f(n)+g(n)\) es \(O(h(n))\)}
\label{sec:orga56e7a1}
\begin{itemize}
\item Si \(f(n)\) es \(O(h(n))\), existen las constantes \(c_f\) y \(n_f\) tal que
\begin{itemize}
\item \(f(n) \leq c_f h(n)\) para toda \(n \geq n_f\)
\end{itemize}
\item Si \(g(n)\) es \(O(h(n))\), existen las constantes \(c_g\) y \(n_g\) tal que
\begin{itemize}
\item \(g(n) \leq c_g h(n)\) para toda \(n \geq n_g\)
\end{itemize}
\end{itemize}

\noindent
Si sumamos ambas desigualdades, se obtiene:
$$
f(n) + g(n) \leq c_f h(n) + c_g h(n) \quad \mbox{para toda} \quad
n \geq \max{n_f, n_g}
$$
que al manipular algebráicamente queda:
$$
f(n) + g(n) \leq (c_f + c_g) h(n) \quad \mbox{para toda} \quad
n \geq \max\{n_f, n_g\}
$$

\noindent
Por definición, se comprueba que \(f(n) + g(n)\) es \(O(h(n))\)
\subsection{Si \(f(n)\) es \(\Omega(h(n))\) y \(g(n)\) es \(\Omega(h(n))\), entonces \(f(n)+g(n)\) es \(\Omega(h(n))\)}
\label{sec:org4fdfb90}
\begin{itemize}
\item Si \(f(n)\) es \(\Omega(h(n))\), existen las constantes \(c_f'\) y \(n_f'\) tal que
\begin{itemize}
\item \(c_f' h(n) \leq f(n)\) para toda \(n \geq n_f'\)
\end{itemize}
\item Si \(g(n)\) es \(\Omega(h(n))\), existen las constantes \(c_g'\) y \(n_g'\) tal que
\begin{itemize}
\item \(c_g' h(n) \leq g(n)\) para toda \(n \geq n_g'\)
\end{itemize}
\end{itemize}

\noindent
Si sumamos ambas desigualdades, obtenemos:
$$
c_f' h(n) + c_g' h(n) \leq f(n) + g(n) \quad \mbox{para toda} \quad
    n \geq \max\{n_f', n_g'\}
$$
que al manipular algebráicamente queda:
$$
(c_f' + c_g') h(n) \leq f(n) + g(n) \quad \mbox{para toda} \quad
    n \geq \max\{n_f', n_g'\}
$$

\noindent
Por definición, se comprueba que \(f(n) + g(n)\) es \(\Omega(h(n))\)
\subsection{Si \(f(n)\) es \(\Theta(h(n))\) y \(g(n)\) es \(\Theta(h(n))\), entonces \(f(n)+g(n)\) es \(\Theta(h(n))\)}
\label{sec:org0ddcad2}
\begin{itemize}
\item Si \(f(n)\) es \(\Theta(h(n))\), existe las constantes \(c_f', c_f\) y \(n_f\) tal que
\begin{itemize}
\item \(c_f' h(n) \leq f(n) \leq c_f h(n)\) para toda \(n \geq n_f\)
\end{itemize}
\item Si \(g(n)\) es \(\Theta(h(n))\), existe las constantes \(c_g', c_g\) y \(n_g\) tal que
\begin{itemize}
\item \(c_g' h(n) \leq g(n) \leq c_g h(n)\) para toda \(n \geq n_g\)
\end{itemize}
\end{itemize}

\noindent
Si sumamos ambas desigualdades tenemos:
$$
c_f' h(n) + c_g' h(n) \leq f(n) + g(n) \leq c_f h(n) + c_g h(n)
    \quad \mbox{para toda} \quad
    n \geq \max\{n_f, n_g\}
$$
si a su vez, esta nueva desigualdad la manipulamos algebraicamente, nos queda
$$
(c_f' + c_g') h(n) \leq f(n) + g(n) \leq (c_f + c_h) h(n)
    \quad \mbox{para toda} \quad
    n \geq \max\{n_f, n_g\}
$$

\noindent
Por definición, se comprueba que \(f(n) + g(n)\) es \(\Theta(h(n))\)



\end{document}
